\section{Sharp entropy interpolation under noise}
\label{flow:section}

The sharp production bound has a direct dynamical meaning: a soft
field loses information at least as fast as a noisy coordinate with
the same information.  The appropriate coordinate for this comparison
is the inverse entropy deficit, rather than information itself.
In that coordinate the comparison flow is simply exponential decay.

\begin{proof}[Proof of Theorem~\ref{main:soft}]
Write $g=u$ for the field in the theorem.
First suppose $0<\rho<1$, and write $\rho=e^{-t}$ with $t>0$.
If $g$ is constant, both sides of \eqref{main:interpolation} are zero.
Otherwise every value of $g_s=T_{e^{-s}}g$ lies in $(-1,1)$ for
$s>0$: the noise kernel gives strictly positive weight to every input,
and a nonconstant $[-1,1]$-valued input cannot average to either
endpoint.  The mean stays fixed.  By
\eqref{main:flow-identity} and Theorem~\ref{upgrade:production},
\begin{equation}
 \frac{d}{ds}I(g_s)=-D(g_s)\le-F(I(g_s))
 \qquad(s>0).
 \label{flow:sharp-trajectory}
\end{equation}

For a nonconstant input, $g_s$ is nonconstant at every finite time,
because the Walsh characters $\chi_A(x)=\prod_{i\in A}x_i$,
$A\subseteq\{1,\ldots,n\}$, form a basis and satisfy
$T_{e^{-s}}\chi_A=e^{-s|A|}\chi_A$. All these eigenvalues are positive.  Strict convexity of $\psi$ therefore gives
$0<I(g_s)<L$.  Set $r(s)=\psi^{-1}(I(g_s))\in(0,1)$.
The identity $F(\psi(r))=rV(r)$ and $\psi'(r)=V(r)$ turn
\eqref{flow:sharp-trajectory} into
\[
 V(r(s))r'(s)\le-r(s)V(r(s)),\qquad r'(s)\le-r(s).
\]
Thus, for $0<\varepsilon<t$,
\[
 r(t)\le e^{-(t-\varepsilon)}r(\varepsilon).
\]
On a finite cube the noise operator and the continuous functional
$I$ converge to their initial values as $\varepsilon\downarrow0$.
Continuity of $\psi^{-1}$, including at $L$, yields
\[
 \psi^{-1}(I(T_{e^{-t}}g))
 \le e^{-t}\psi^{-1}(I(g)).
\]
Applying the increasing function $\psi$ proves the claim for
$0<\rho<1$.  This positive-time argument does not require the
initial entropy production to be finite.

At $\rho=0$, the field is constant; at $\rho=1$, the assertion
is an identity.  For negative correlation,
$T_{-\rho}g(x)=T_\rho g(-x)$, and the change of variables
$x\mapsto-x$ preserves the uniform measure and information.
This proves the full stated range.

For $g(x)=a x_i$ with $|a|\le1$, the initial information is
$\psi(|a|)$ and the information after noise is
$\psi(|\rho a|)$. Thus every soft coordinate attains equality.
\end{proof}

\begin{proof}[Proof of Corollary~\ref{main:ck}]
Since $f$ is Boolean, $H(f(x))=0$ at every input, so
$I(f)=H(m)$.  Also $I(f(X);Y)=I(T_\rho f)$.  The first
inequality is therefore Theorem~\ref{main:soft}.
The second follows from $r_m\le1$ and monotonicity of $\psi$.
If $0<|m|<1$, then $0<H(m)<L$ and $0<r_m<1$, giving the
stated strict comparison with the unbiased bound.
For a signed coordinate $f$, $T_\rho f=\rho f$, so its
information is exactly $\psi(|\rho|)$.
\end{proof}

\subsection{Comparison of mean-dependent production bounds}
\label{flow:comparison}

Theorem~\ref{upgrade:production} gives more than the scalar rate needed
for Theorem~\ref{main:soft}. For $0<m<1$ and $0\le I<H(m)$, put
$h=H(m)-I$. The hybrid production profile of
Chen et al.\ \cite[Section~9 and Lemma~9.2]{ChenEtAl2026CK}, converted
to natural logarithms and the generator $N$ used here, is
\[
 C_m(I)=\max\{F(L-h)-mV(s),F(I)\},\qquad mH(s)=hs,
 \qquad m\le s<1.
\]
The equation for $s$ has a unique solution: $H(s)/s$ decreases
strictly from $H(m)/m$ to zero on $[m,1)$. In particular, the
$F(I)$ branch yields Theorem~\ref{main:soft} by the inverse-entropy
integration above, with the initial information left unchanged.

The mean-dependent estimate in Theorem~\ref{upgrade:production}
has a different strength. At every fixed $m\in(0,1)$,
\begin{equation}
 \mathcal B_m(H(m)-h)\sim\frac{1-m^2}{2}\log\frac1h,
 \qquad C_m(H(m)-h)\sim\frac{1-m}{2}\log\frac1h
 \quad(h\downarrow0).
 \label{flow:profile-asymptotics}
\end{equation}
Here is a direct verification. For $r=1-\varepsilon$,
\[
 H(r)=\frac\varepsilon2\log\frac{2\exp(1)}{\varepsilon}
       +O(\varepsilon^2),\qquad
 V(r)=\frac12\log\frac2\varepsilon+O(\varepsilon).
\]
Consequently $F(L-h)\sim\tfrac12\log(1/h)$.
The equation $H(s)=hs/m$ gives $s\to1$ and
$V(s)\sim\tfrac12\log(1/h)$. As $F(H(m)-h)$ remains bounded,
these facts give the second equivalent in
\eqref{flow:profile-asymptotics}. The first follows from
\[
 \mathcal B_m(H(m)-h)
 =2(H(m)-h)+(1-m^2)S\left(L-\frac h{1-m^2}\right)
\]
and $S(u)=F(u)-2u$ for positive $u$.
Thus the displayed profile $C_m$ does not dominate
$\mathcal B_m$: the ratio $\mathcal B_m(H(m)-h)/C_m(H(m)-h)$
tends to $1+m>1$ as $h\downarrow0$.
This comparison concerns the two explicit production profiles.
